\documentclass[12pt]{article}
\usepackage{graphicx} % Required for inserting images
\usepackage[a4paper,margin=1in]{geometry}
\usepackage{multicol}
\usepackage{amsmath}
\usepackage{fancyhdr}
\usepackage{lastpage}  % Ensures total page count

\title{5730 - Mobile Communications and Internet of Things Spring Exam 2025}
\date{18. March 2025}

\begin{document}
\pagestyle{fancy}
\fancyhf{}  % Clear default headers and footers
\renewcommand{\headrulewidth}{0pt}  % No top horizontal line
\renewcommand{\footrulewidth}{0pt}  % No bottom horizontal line
\fancyfoot[R]{\thepage/\pageref{LastPage}}  % Right footer: X/Y format
\fancypagestyle{plain}{
    \fancyhf{}  % Clear default first-page footer
    \fancyfoot[R]{\thepage/\pageref{LastPage}}  % Apply the same footer
    \renewcommand{\headrulewidth}{0pt}  % Remove top line on first page
    \renewcommand{\footrulewidth}{0pt}  % Ensure no bottom line
}
\maketitle

\section*{Instructions to Students}

\begin{itemize}
    \item Ensure clarity and precision in your answers.
    \item Attempt all sections. Each of the three sections contributes one-third (1/3) of the total mark.
    \item All answers must be written in \textbf{English}.
    \item Write your answers in the provided answer sheets using a pen.
    \item Each page must contain your name and student number.
    \item Number all submitted pages as \textit{Page X/Y}, where \textit{X} is the page number and \textit{Y} is the total (e.g., \textit{1/9, 2/9, ..., 9/9} for nine pages). 
    \item For the multiple-choice section, indicate your answers using the format: Q1 A, Q2 C, ...
    \item In the appendix of this document, you will find a logarithm tables for reference.
\end{itemize}

\newpage
\section{Multiple Choice Questions (MCQs)}
\begin{multicols}{2}
\begin{enumerate}
    \item Which of the following is NOT a key feature of 5G networks?
    \begin{enumerate}
        \item eMBB
        \item URLLC
        \item mMTC
        \item IPv4
    \end{enumerate}
    
    \item Which of the following is a key reason why 5G networks support both standalone (SA) and non-standalone (NSA) deployment modes?
    \begin{enumerate}
        \item To allow gradual migration from LTE to 5G while maintaining service continuity.
        \item To enable devices to switch dynamically between 4G and 5G spectrum in real time.
        \item To ensure that all 5G devices operate exclusively in mmWave frequencies.
        \item To reduce the need for small cell deployments in urban environments.
    \end{enumerate}

    \item What is the primary function of network slicing in 5G?
    \begin{enumerate}
        \item Improving battery life by dynamically adjusting power consumption based on real-time network traffic conditions
        \item Enabling multiple virtual networks on a single physical infrastructure
        \item Encrypting data using end-to-end security protocols that prevent unauthorized access across network slices.
        \item Optimizing spectrum usage by dynamically reallocating frequency resources to devices based on current traffic demand.
    \end{enumerate}

    \item Which IoT protocol in licenced spectrum is specifically optimized for low-power wide-area networks?
    \begin{enumerate}
        \item NB-IoT
        \item HTTP
        \item Zigbee
        \item LoRa
    \end{enumerate}

    \item What is the main purpose of beamforming in 5G networks?
    \begin{enumerate}
        \item Data compression
        \item Increasing antenna gain
        \item Encrypting data packets
        \item Load balancing across devices
    \end{enumerate}
    
    \item Which of the following is a key feature of edge computing?
    \begin{enumerate}
        \item Centralized processing
        \item Low-latency data processing close to the source
        \item Long-term data storage
        \item Real-time encryption
    \end{enumerate}

    \item What is the typical range of NB-IoT in urban areas?
    \begin{enumerate}
        \item 10-100 meters
        \item 1-10 kilometers
        \item 50-100 kilometers
        \item 100+ kilometers
    \end{enumerate}

    \item URLLC in 5G is ideal for:
    \begin{enumerate}
        \item High-speed file downloads
        \item Autonomous driving
        \item Smart metering
        \item Remote learning
    \end{enumerate}
    \newpage

    \item AI/ML can optimize IoT networks by:
    \begin{enumerate}
        \item Predicting network congestion and optimizing routing  
        \item Decreasing encryption standards
        \item Simplifying device manufacturing by automating the design and testing of IoT hardware.
        \item Replacing all physical IoT sensors with virtual models that simulate real-world conditions.
    \end{enumerate}

    \item Quantum communication in cellular networks primarily aims to:
    \begin{enumerate}
        \item Increase bandwidth
        \item Improve security
        \item Enable faster video streaming
        \item Reduce hardware costs
    \end{enumerate}

    \item Massive MIMO improves network performance by:
    \begin{enumerate}
        \item Allowing simultaneous transmission to multiple users
        \item Increasing spectrum availability by dynamically reallocating unused frequency bands.
        \item Improving signal to noise ratio to a individual device by using multiple antennas
        \item Encrypting data at the physical layer to enhance security in wireless transmission.   
    \end{enumerate}
    \item Data analytics in IoT can provide real-time insights by:
    \begin{enumerate}
        \item Processing data on centralized servers
        \item Using edge computing for local processing
        \item Encrypting all network traffic
        \item Storing data indefinitely
    \end{enumerate}
    \item What is the primary focus of sustainability in future cellular networks?
    \begin{enumerate}
        \item Increasing the number of devices
        \item Reducing energy consumption
        \item Simplifying device design
        \item Enabling faster network slicing
    \end{enumerate}

    \item Which technology enables dynamic resource allocation in 5G?
    \begin{enumerate}
        \item Beamforming
        \item SDN
        \item Blockchain
        \item IPv6
    \end{enumerate}

    \item NB-IoT is best suited for which type of applications?
    \begin{enumerate}
        \item High-bandwidth video streaming
        \item Low-power, long-range communication
        \item Real-time gaming
        \item Blockchain-based networks
    \end{enumerate}

    \item What is the role of eMBB in 5G?
    \begin{enumerate}
        \item Providing high data rates for enhanced multimedia experiences
        \item Supporting low-power devices
        \item Enabling ultra-low latency
        \item Managing security protocols
    \end{enumerate}

    \item What does MQTT stand for?
    \begin{enumerate}
        \item Multi Query Telemetry Transport
        \item Message Queue Telemetry Transport
        \item Message Quick Transfer Technology
        \item Multi Quick Telemetry Transport
    \end{enumerate}

    \newpage
    \item How does AI/ML contribute to network optimization in IoT?
    \begin{enumerate}
        \item By reducing the need for encryption
        \item By enabling predictive maintenance and efficient routing
        \item By simplifying device designs
        \item By eliminating the need for data processing
    \end{enumerate}

    \item What is the key advantage of URLLC in 5G?
    \begin{enumerate}
        \item High data rates
        \item Support for massive device connectivity
        \item Ultra-low latency and high reliability
        \item Energy efficiency
    \end{enumerate}
    
    \item Which component in MQTT is responsible for managing message distribution between publishers and subscribers? 
    \begin{enumerate}
    \item Client
    \item Broker
    \item Gateway
    \item Proxy server
    \end{enumerate}

    \item Which factor affects the range of wireless communication the most?
\begin{enumerate}
\item Packet size 
\item Frequency attenuation and penetration losses 
\item Network topology 
\item Transmission protocol 
\end{enumerate}

    \item What is the primary advantage of NB-IoT over LoRaWAN and Sigfox?
\begin{enumerate}
\item It uses licensed spectrum for improved reliability. 
\item It supports unlimited devices per cell.
\item It has the highest data rate in LPWAN technologies.
\item It does not require a base station for communication.
\end{enumerate}
\item According to Shannon’s Capacity Theorem, what factors influence the maximum achievable data rate in a wireless communication system?
\begin{enumerate}
    \item Bandwidth and signal-to-noise ratio (SNR)
    \item Modulation scheme and power amplifier type
    \item The number of users in the network
    \item The type of network slicing used
\end{enumerate}

\item In the context of blockchain-based authentication for IoT devices, which of the following is the most significant advantage compared to traditional centralized Public Key Infrastructure?
\begin{enumerate}
    \item Blockchain completely eliminates the need for cryptographic key exchange between devices
    \item Blockchain allows for immutable and decentralized identity verification without a single point of failure
    \item Blockchain-based authentication reduces computational load by avoiding any use of digital signatures
    \item Blockchain enforces dynamic IP address allocation, enhancing identity privacy in IoT environments.
\end{enumerate}
\newpage
\item Which of the following statements best describes a primary trade-off involved in deploying lightweight symmetric encryption for IoT security in 5G networks?
\begin{enumerate}
    \item Increased encryption strength leads to lower energy consumption, making these ciphers ideal for ultra-low-power devices.
    \item Lightweight ciphers offer better forward secrecy than standard block ciphers like AES.
    \item Lower computational and memory overhead comes at the cost of reduced cryptographic strength or resistance to certain attacks.
    \item These ciphers are incompatible with the high-throughput requirements of 5G URLLC scenarios.
\end{enumerate}
\end{enumerate}
\end{multicols}

\section{Short Questions}
\begin{enumerate}
    \item Briefly explain the role of eMBB in 5G communication.
    \item Define edge computing and its significance in IoT applications.
    \item What is the purpose of beamforming in massive MIMO?
    \item Mention two real-world applications of massive IoT.
    \item What is the relationship between bandwidth and data rate in a wireless communication system?
    \item Compare QPSK and 16-QAM in terms of spectral efficiency and robustness to noise. 
    \item How does massive MIMO contribute to energy efficiency in 5G networks?
    \item What is the difference between SDN and traditional network architecture?
    \item What are the advantages and disadvantages of using licensed vs. unlicensed spectrum for IoT applications? Provide one example for each.
    \item SHA-256 is widely used for cryptographic hashing and authentication in IoT security. Explain its role in authentication and discuss its limitations in IoT applications.
\end{enumerate} 
\newpage
\section{Logical Reasoning}
Consider an IoT parking sensor that must communicate with a 5G base station located 1 km away.
The sensor has a transmission power of 20 dBm, and the base station transmits at 23 dBm. The antenna gain of the sensor and base station is 0 dBi. The carrier frequency is 3 GHz. Various losses in the systems amount to 8 dB.
\begin{enumerate}
    \item The IoT device has a receive bandwidth of 180 kHz. Calculate the noise floor at room temperature (290K).
    \item Using the same IoT sensor and base station setup:
    \begin{enumerate}
        \item Compute the Free-Space Path Loss (FSPL). For ease of calculation you can assume that the speed of light $c \approx 3\cdot10^8$ and $\pi \approx 3$.
        \item If the transmitting device is aiming for an SNR of 0 dB at the receiver, calculate the required transmission power $P_t$.
        \item Explain what happens to the SNR and data rate if the bandwidth is increased to 1.4 MHz, provided that the transmission power remains as calculated in 2b. Does increasing bandwidth always improve the data rate?
    \end{enumerate}
    \item Assume the system supporting the IoT network uses 16-QAM modulation.
    \begin{enumerate}
        \item How many bits per symbol does this modulation scheme support?
        \item Would increasing the modulation order to 64-QAM improve or worsen the communication for the IoT sensor? Explain.
    \end{enumerate}
\item A 5G system has a received power that results in an SNR of 0 dB. However, the system requires an SNR of 10 dB for reliable communication.
    \begin{enumerate}
        \item Compute the number of beamforming antennas required to achieve the required SNR.
        \item Assuming a bandwidth of 100 MHz, calculate the theoretical channel capacity for SNR = 10 dB.
    \end{enumerate}

    \end{enumerate}
\newpage
\section{Supplemental material}
\subsection{Logarithm Properties}
\begin{itemize}
    \item \textbf{Product Rule:} 
    \[
    \log_b (xy) = \log_b (x) + \log_b (y)
    \]
    \item \textbf{Quotient Rule:} 
    \[
    \log_b \left( \frac{x}{y} \right) = \log_b (x) - \log_b (y)
    \]
    \item \textbf{Power Rule:} 
    \[
    \log_b (x^n) = n \log_b (x)
    \]
    \item \textbf{Base Change Formula:} 
    \[
    \log_b (x) = \frac{\log_c (x)}{\log_c (b)}
    \]
    \item \textbf{Log of Reciprocal:} 
    \[
    \log_b \left( \frac{1}{x} \right) = -\log_b (x)
    \]
    \item \textbf{Logarithm to Exponential Conversion:}
    \[
    x \cdot \log_b(y) = z \quad \Rightarrow \quad y = b^{\frac{z}{x}}
    \]
\end{itemize}
\subsection{Commonly Used Constants and Units}

\begin{itemize}
    \item \textbf{Speed of Light:} 
    \[
    c \approx 3.0 \times 10^8 \text{ m/s}
    \]
    \item \textbf{Boltzmann's Constant:} 
    \[
    k_B = 1.38 \times 10^{-23} \text{ J/K}
    \]
        \item \textbf{Thermal Noise Power Density:}  
    \[
    P_n (1 \text{ Hz}) = -174 \text{ dBm/Hz}
    \]
    
\end{itemize}

\subsection{Logarithm Conversion Examples}
\begin{itemize}
    \item \(\log_{2} (16) = 4\), since \(2^4 = 16\).
    \item \(\log_{10}(12) = \log_{10}(6)+\log_{10}(2)\)
    \item Converting \( 20 \) dB to a linear factor: \(10^{\frac{20}{10}} = 10^{2} = 100\)

\end{itemize}

\begin{table}[h]
    \centering
    \renewcommand{\arraystretch}{1.2}
    \begin{tabular}{|l|r|r|r|r|}
        \hline
        \textbf{Value} & \(\log_{10}(\text{Value})\) & \(\log_{2}(\text{Value})\) & Scientific Notation & SI - Prefix \\
        \hline
        \(10^{-9}\)   & -9.000  & -29.897 & \(1.0 \times 10^{-9}\) & Nano \\
        \(10^{-8}\)   & -8.000  & -26.575 & \(1.0 \times 10^{-8}\) & \\
        \(10^{-7}\)   & -7.000  & -23.253 & \(1.0 \times 10^{-7}\) & \\
        \(10^{-6}\)   & -6.000  & -19.931 & \(1.0 \times 10^{-6}\) & Micro \\
        \(10^{-5}\)   & -5.000  & -16.609 & \(1.0 \times 10^{-5}\) & \\
        \(10^{-4}\)   & -4.000  & -13.287 & \(1.0 \times 10^{-4}\) & \\
        \(10^{-3}\)   & -3.000  & -9.966  & \(1.0 \times 10^{-3}\) & Milli \\
        \(0.01\)   & -2.000  & -6.644  & \(1.0 \times 10^{-2}\) & \\
        \(0.1\)   & -1.000  & -3.322  & \(1.0 \times 10^{-1}\) & \\
        \hline
        1             & 0.000   & 0.000   & \(1.0 \times 10^{0}\) & - \\
        2             & 0.301   & 1.000   & \(2.0 \times 10^{0}\) & \\
        3             & 0.477   & 1.585   & \(3.0 \times 10^{0}\) & \\
        4             & 0.602   & 2.000   & \(4.0 \times 10^{0}\) & \\
        5             & 0.699   & 2.322   & \(5.0 \times 10^{0}\) & \\
        6             & 0.778   & 2.585   & \(6.0 \times 10^{0}\) & \\
        7             & 0.845   & 2.807   & \(7.0 \times 10^{0}\) & \\
        8             & 0.903   & 3.000   & \(8.0 \times 10^{0}\) & \\
        9             & 0.954   & 3.170   & \(9.0 \times 10^{0}\) & \\
        10            & 1.000   & 3.322   & \(1.0 \times 10^{1}\) & \\
        \hline
        \(100\)      & 2.000   & 6.644   & \(1.0 \times 10^{2}\) & \\
        \(1000\)      & 3.000   & 9.966   & \(1.0 \times 10^{3}\) & Kilo \\
        \(10^4\)      & 4.000   & 13.287  & \(1.0 \times 10^{4}\) & \\
        \(10^5\)      & 5.000   & 16.609  & \(1.0 \times 10^{5}\) & \\
        \(10^6\)      & 6.000   & 19.931  & \(1.0 \times 10^{6}\) & Mega \\
        \(10^7\)      & 7.000   & 23.253  & \(1.0 \times 10^{7}\) & \\
        \(10^8\)      & 8.000   & 26.575  & \(1.0 \times 10^{8}\) & \\
        \(10^9\)      & 9.000   & 29.897  & \(1.0 \times 10^{9}\) & Giga \\
        \hline
    \end{tabular}
\end{table}

\end{document}
